\vspace{-0.5em}
\section{Evaluation and Results}\label{sec:evaluation}

Our evaluation shows that PaTAS produces trust estimates that converge during training, respect the properties of \cref{sec:ptastheory}, and remain interpretable under noisy features, corrupted labels, and adversarial perturbations. We vary input trustworthiness from fully trusted through fully uncertain to fully distrusted, track the trust, uncertainty, and distrust mass, and compare against post-training accuracy.

We conduct four experiments of increasing complexity:
\begin{enumerate}
  \item \emph{Breast Cancer Classification}, to assess behavior on a small, tabular medical dataset.
  \item \emph{MNIST Digit Classification}, to evaluate PaTAS across multiple neural architectures under controlled uncertainty.
  % and randomized trust assignments.
  \item \emph{Poisoned MNIST}, to evaluate PaTAS and IPTA in the presence of adversarial corruption and data poisoning.
  \item \emph{Comparative Trust Evaluation}, to quantify PaTAS against established uncertainty baselines on out-of-distribution and corruption detection, and to characterize when input-level trust is informative.
\end{enumerate}

\vspace{-0.5em}
\subsection{Experimental Setup }
\subsubsection{Experiment 1 - Breast Cancer Classification}\label{exp:1}
We use the \href{https://archive.ics.uci.edu/dataset/17/breast+cancer+wisconsin+diagnostic}{Breast Cancer Wisconsin (Diagnostic) Dataset}~\cite{breastcancer-wisconsin}, containing 569 samples with 30 numeric features (e.g., radius, area, symmetry) derived from breast mass fine needle aspirates. The neural network classifies the tumors into benign or malignant categories.
The neural network architecture consists of 30 input neurons, 16 hidden neurons, and 2 output neurons, with ReLU activation in the hidden layer and Softmax in the output. The model is trained for 15 epochs with a batch size of 64 and a learning rate of 0.2, achieving 97\% test accuracy when the data are not modified.

We degrade the training data in controlled ways and assign corresponding trust assessments, combining three extreme profiles for features and labels alike, fully trusted \((1,0,0)\), fully distrusted \((0,1,0)\), and fully uncertain \((0,0,1)\), into nine cases. These dogmatic and vacuous opinions are the canonical boundary cases in SL. Two intermediate scenarios are added below to show behavior under partial trust.

Feature and label degradation is simulated with controlled perturbation functions: uncertain features receive additive uniform noise with clipping to the valid feature range, distrusted features are resampled uniformly over the feature space, and labels are degraded analogously through random noise (uncertainty) or complete replacement (distrust). The exact perturbation functions are given in Appendix~\ref{app:perturb}.

Finally, to complement the boundary cases, we evaluate two intermediate trust scenarios: 
\begin{enumerate}[label=\roman*.]
    \item the same scenario for a fully uncertain input features and fully uncertain labels, but where the assessment is set to \((0.25, 0.25, 0.5)\), reflecting partial distrust and trust, instead of a fully uncertain opinion (0,0,1) and,
    \item a mild degradation, where features are perturbed with probability 0.15 and the trust assessment is set to \((0.25, 0, 0.75)\).
    
\end{enumerate}

\subsubsection{Experiment 2 - MNIST}\label{exp:2}

We evaluate PaTAS on MNIST~\cite{lecun1998mnist} (60{,}000 training and 10{,}000 test images of \(28\times28\) pixels, ten digit classes) using four architectures of increasing width: 784-16-10, 784-32-10, 784-64-10, and 784-128-10. Although not state of the art for MNIST, they suffice to expose how trust propagates across model sizes. Each uses ReLU in the hidden layer and Softmax at the output, trained for 20 epochs with batch size 128 and learning rate 0.05, reaching 99\% test accuracy (when training data are fully trusted); the gradient-evidence threshold is $\epsilon = 0.05$ here and in Experiments 3 and 4. Trust assessments for both features and labels are fully uncertain, corresponding to no knowledge of the dataset; for contrast we repeat the smallest architecture with a fully trusted assessment.

\subsubsection{Experiment 3 - Poisoned MNIST}\label{exp:3}
We poison one third of the MNIST training images by flipping the labels of digits 6 and 9 and simultaneously adding a visible patch of fixed size at the top-left corner, following common data-poisoning and backdoor practice~\cite{chen2017targetedbackdoorattacksdeep}; the remaining two thirds stay clean. This lets us examine how PaTAS responds to corrupted labels and adversarial triggers. We use the 784-128-10 architecture. Patch pixels and the labels of patched 6s and 9s are assigned distrust, everything else trust, so that PaTAS focuses on the corrupted regions while trusting the unaffected data.


\subsubsection{Experiment 4 - Comparative Trust Evaluation}\label{exp:4}
We compare the per-sample PaTAS assessment against three established uncertainty baselines: the maximum softmax probability, MC~Dropout~\cite{gal2016dropoutbayesianapproximationrepresenting} ($T{=}30$ passes, $p{=}0.2$), and Evidential Deep Learning~\cite{sensoy2018evidentialdeeplearningquantify}. Each baseline trains its own network with the modifications its method requires (dropout active during training for MC~Dropout, the evidential Dirichlet loss for EDL) under the shared architecture, optimizer, and epoch budget. DeepTrust~\cite{10.3389/frai.2020.00054} has no public implementation and is discussed in \cref{sec:rw}. Primary datasets are MNIST and Fashion-MNIST~\cite{xiao2017fashionmnist} (784-512-10, 20 epochs); GTSRB~\cite{GTSRB} (1024-128-43, features standardized with train statistics) serves as an applicability-boundary case discussed below.

The PaTAS assessment instantiates route~(ii), estimation from data, of \cref{sec:acquisition}: each input feature receives a trust opinion from its conformity to the training-data statistics, mapped through the same Baseline-Prior Quantification as parameter trust (\cref{eq:q1}), with per-feature evidence weighted by the feature's reliability as a witness, $w_j = (\sigma_{\mathrm{ref}}/\sigma_j)^2$ (inverse-variance weighting with $\sigma_{\mathrm{ref}}$ the median effective deviation, capped at 64; the full construction and all parameters are specified in \cref{app:conformity}). The trust factor is the serial discount $P = P_{\mathrm{model}} \cdot P_{\mathrm{input}}$ of the IPTA path trust under a fully trusted input and the pooled input conformity, and the reported score anchors the discounted confidence at the class prior, $P \cdot \hat{p} + (1-P)/K$. Every method is evaluated identically on clean test data, under Bernoulli-uniform feature corruption ($p \in \{0.3, 0.6\}$, amplitude scaled to the feature range), on a mixed batch pooling clean and corrupted inputs 1:1, and as a detector of out-of-distribution inputs (the OOD set passes through the identical preprocessing pipeline: Fashion-MNIST for MNIST and vice versa); post-hoc isotonic calibration is fit for each method on a held-out clean split. Results are mean~$\pm$~std over three evaluation seeds.

\paragraph*{Implementation Details for All Experiments}
During inference, multiplication operations in the feedforward phase are implemented using SL trust discounting, while addition operations employ a generalized SL averaging fusion operator to support summations over multiple inputs (see \cref{def:trustfunc}). Trust revision uses the same averaging fusion. The PaTAS framework and Neural Network were implemented in Python 3.13 using NumPy. Two main modules were developed: \texttt{PrimaryNN}, which handles network structure, training, and inference, and \texttt{PaTAS}, which implements trust assessment and propagation functions defined in the operational flow (\cref{fig:ptasoperation}). These modules interact to dynamically evaluate and update trust during feedforward and backpropagation. The full implementation, including all modules, dependencies, and experiment scripts, is available as \suppp, with detailed instructions for reproducing all experiments.

\vspace{-0.5em}
\subsection{Results and Analysis}
Evaluation is performed during training: after each feedforward/backpropagation iteration we assess the trustworthiness of the resulting output under three boundary input profiles, fully trusted, fully uncertain, and fully distrusted, which are the most informative extreme cases. For each, we track the trust, uncertainty, and distrust mass over training, giving nine curves per experiment (Experiments 1--3)\footnote{The base rate is set to $0.5$ and held constant across all experiments.}. Results are collected in the \supp and summarized in \cref{tab:cancer_results,tab:mnist,tab:mnistpois,tab:mnistpoisipta}; they confirm several properties proven in \cref{sec:ptastheory}:

\begin{table*}[ht]
\begin{minipage}{0.63\textwidth}
\centering
\begin{tabular}{cccccc}
\hline
X Trust Assessment & Y Trust Assessment& \(\epsilon = 0.01\) & \(\epsilon = 0.1\) & Train (\%) & Test (\%) \\
\hline
fully distrusted & any & 0 & 0 & 47--58 & 41--55 \\
any & fully distrusted & 0 & 0 & 96--98 & 3--4 \\
fully uncertain & fully uncertain & 0.17 & 0.28 & 66 & 85 \\ % noise-noise
% \hline
fully uncertain & fully trusted & 0.15 & 0.32 & 96 & 96 \\ % noise-clean
% \hline
 fully trusted & fully uncertain & 0.12 & 0.29 & 69 & 94 \\ % clean-noise
% \hline
fully trusted & fully trusted & 0.62 & 0.88 & 98 & 97 \\ % clean-clean
(0.25, 0.25, 0.5) & (0.25, 0.25, 0.5)& 0.06 & 0.11 & 66 & 85 \\
% \hline
(0.25, 0, 0.75) & (0.25, 0, 0.75)& 0.20 & 0.35 & 75 & 94\\
% \hline

\hline
\end{tabular}
\caption{Final trust mass and train/test accuracy (\%) for the Breast Cancer experiment. Rows where either assessment is fully distrusted collapse to zero trust and are grouped.}
\label{tab:cancer_results}
\end{minipage}
\hfill
\begin{minipage}{0.37\textwidth}
\centering
\begin{tabular}{cccc}
\hline
Hidden Neurons & Trust \(t\) & Train (\%) & Test (\%) \\
\hline
16 & 0.306 & 65 & 72\\
% \hline
32 & 0.325 & 66 & 89 \\
% \hline
64 & 0.333 & 67 & 90\\
% \hline
128 & 0.337 & 68 & 93\\
% \hline
16 (trust) & 0.778 & 96 & 95\\
\hline
\end{tabular}
\caption{Summary of final trust mass for the MNIST Classification Experiment.}
\label{tab:mnist}
\end{minipage}

\end{table*}

\begin{itemize}
\item \emph{Convergence of Trust Assessment}: when accuracy converges, trust assessment converges as expected.
\item \emph{Inference on fully uncertain input yields fully uncertain output}: a fully uncertain input always produces a fully uncertain output.
\item \emph{Symmetric Inference}: when input trust assessments are symmetric, the inference results remain symmetric.
\end{itemize}

We therefore report only the fully trusted input case in detail: uncertain inputs yield uncertain outputs and symmetry makes the distrusted case a mirror of the trusted one (\cref{thm:infvac}), and under trusted inputs the distrust mass stays close to zero, so the trust mass alone determines the opinion. \Cref{tab:cancer_results,tab:mnist,tab:mnistpois,tab:mnistpoisipta} report the converged trust mass.


\subsubsection*{Experiment 1 -- Breast Cancer Classification}
Corrupted labels mislead rather than prevent learning: training accuracy stays high (96--98\%) while test accuracy collapses to near zero. Corrupted features eliminate learning in every setting (41--55\% test, near chance for the binary task), while label noise depresses training accuracy (69\%) yet still permits generalization (94\%).

We evaluated the system for two values of $\epsilon$ (0.1 and 0.01), as summarized in \cref{tab:cancer_results}. We recall that $\epsilon$ is the threshold used by the NodeTrust function in \cref{alg:trustupdate}.
\begin{itemize}
\item For $\epsilon = 0.1$: when $X$ is uncertain, trust mass stabilizes around 0.28 for uncertain $Y$ and 0.32 for trusted $Y$. Fully trusted $X$ gives 0.29 for uncertain $Y$ and $0.88$ for trusted $Y$. Whatever the $\epsilon$ value, when $X$ or $Y$ are distrusted the trust mass rapidly falls to 0.
\item For $\epsilon = 0.01$: the system reads more gradients as large, and all partial-trust configurations drop (0.12--0.17); notably, under degraded features the trusted-$Y$ case falls \emph{below} the uncertain-$Y$ case (0.15 vs.\ 0.17), since trusted labels turn persistent gradient corrections into negative evidence, the mechanism analyzed in \cref{sec:acquisition}.
\end{itemize}
Overall, there is a strong positive correlation between trust mass and test accuracy among the non-distrusted configurations.

Two findings are worth emphasizing. First, \emph{label trust matters more than feature trust}: setting $X$ as fully uncertain while keeping $Y$ trusted yields a higher trust mass (0.32) than the opposite case, where $X$ is trusted but $Y$ uncertain (0.29). Second, \emph{distrust is more damaging than uncertainty}: replacing a fully uncertain opinion $(0,0,1)$ for both $X$ and $Y$ with a mixed assessment $(0.25,0.25,0.5)$ reduces the trust mass from $0.28$ to $0.11$, showing that even partial distrust degrades trust more strongly than uncertainty alone.




\subsubsection*{Experiment 2 -- MNIST Dataset}
In this experiment, PaTAS produces ten output trust opinions, one for each class. However they are almost the same so we record only the aggregated trust mass in the table. As shown in \cref{tab:mnist}, when features and labels are noisy, accuracy improves as model size increases. Test accuracy is consistently greater than training accuracy across all architectures, suggesting training noise makes the model underestimate its learning progress. While trust increases with model size, the gains become progressively smaller (0.306 $\to$ 0.325 $\to$ 0.333 $\to$ 0.337). This contrasts clearly with the fully trusted evaluation on the smallest architecture, which achieves a trust mass of 0.778, far exceeding any uncertain-data configuration regardless of model size. This highlights that \emph{training a smaller architecture on trusted data is significantly more beneficial for trust than training a larger architecture on uncertain data}. Moreover, we can see that larger architectures yield more stable trust assessments, likely due
to smaller average gradient magnitudes.


\subsubsection*{Experiment 3 – Poisoned MNIST Dataset}
As in the previous experiment, PaTAS has ten output trust opinions. These values are informative because the reliability of the inference path vary across labels, especially when some are patched and others are not. In \cref{tab:mnistpois}, we report the trust mass associated with label 3 (a clean class) and label 6 (a poisoned class; see \cref{exp:3}). Alongside overall train and test accuracy, we also report accuracy on clean test samples of digits 3 and 6, as well as accuracy on poisoned versions of those test samples.

The results suggest that PaTAS is able to reflect differences in reliability between the clean and poisoned classes. Across all patch sizes, the trust mass assigned to label 3 consistently remains higher than that assigned to label 6 by a wide margin (e.g., 0.90 vs.\ 0.58 for the $4\times4$ patch), confirming that PaTAS reliably distinguishes clean from poisoned inference paths. As shown in \cref{tab:mnistpois}, both trust values decrease as the patch size grows, since larger patches affect a greater number of neurons along the inference path, thereby reducing the overall reliability of the parameters associated with those paths and progressively affecting even the neurons involved in predicting clean classes; the degradation is largest at $27 \times 27$, where the patch dominates nearly the entire input (0.70 and 0.49 for labels 3 and 6), while the clean--poisoned separation remains interpretable at every size.

To evaluate the IPTA, \Cref{tab:mnistpoisipta} reports both accuracy and the corresponding trust opinions for clean digits (3 and 6) and patched digit 6 datasets. The trust opinion is obtained by feedforwarding a fully trusted opinion through the IPTA derived from clean digit 3, clean digit 6, and patched digit 6 images. The backdoor is near-total in this run (patched 6s are misclassified in 99\% of cases), and the trust assessment exposes it in the sharpest possible way: clean digit-6 samples retain high accuracy (97.4\%) yet carry a markedly lower trust mass ($0.570$) than clean digit-3 samples ($0.886$), because the parameters along the digit-6 inference path were learned from poisoned data: an accurate prediction whose supporting evidence PaTAS refuses to endorse, precisely the reliability gap that confidence scores do not reveal. Patched samples inherit the same reduced trust and elevated uncertainty ($0.571$, $u = 0.429$), and explicitly distrusting the patch pixels before feedforwarding through the IPTA additionally surfaces a non-zero distrust mass ($0.560, 0.011, 0.429$). These results show that PaTAS provides interpretable warnings about poisoned outputs even when accuracy alone looks healthy.

\begin{table}[ht]
\centering
\setlength{\tabcolsep}{4pt}
\begin{tabular}{ccccccc}
\hline
Patch & \multicolumn{2}{c}{Trust} & \multicolumn{2}{c}{Clean acc.\ (\%)} & \multicolumn{2}{c}{Patched acc.\ (\%)} \\
size & 3 & 6 & 3 & 6 & 3 & 6 \\
\hline
\((1 \times 1)\) & 0.897 & 0.577 & 97.92 & 97.29 & 78.22 & 0.84 \\
\((4 \times 4)\) & 0.895 & 0.576 & 97.82 & 97.39 & 39.60 & 1.04 \\
\((10 \times 10)\) & 0.881 & 0.570 & 97.92 & 97.39 & 22.08 & 0.84 \\
\((27 \times 27)\) & 0.697 & 0.490 & 97.92 & 97.29 & 0 & 0 \\
\hline
\end{tabular}
\caption{Poisoned MNIST. Overall train/test accuracy stays in the range 96.4--99.6\%/97.7--97.8\% across all patch sizes and is omitted.}
\label{tab:mnistpois}
\end{table}

\begin{table}[ht]
\centering
\begin{tabular}{ccccccccc}
\hline
 & Accuracy (\%) & Trust  & Distrust & Uncertainty \\
\hline
Clean 3 & 97.82 & 0.886 & 0.0 & 0.114\\

Clean 6  & 97.39 & 0.570 & 0.0 & 0.430\\

6 with patch& 1.04 & 0.571 & 0.0 & 0.429\\

6 with patch & -- & 0.560 & 0.011 & 0.429\\
(patch distrusted) &&&&\\
\hline
\end{tabular}
\caption{IPTA results for a neural network trained on the poisoned dataset with patch size $4\times 4$ (single training run; the backdoor drives patched-6 accuracy to 1\%). Results are obtained by feedforwarding a fully trusted opinion through the IPTA, except for the last row where patch pixels are explicitly distrusted. The poisoned class carries reduced trust even on clean, correctly classified inputs.}
\label{tab:mnistpoisipta}
\end{table}

\subsubsection*{Experiment 4 -- Comparative Trust Evaluation}
\Cref{tab:comparative} reports the two detection tasks in which trust and confidence should differ: recognizing out-of-distribution inputs, and recognizing corrupted inputs inside a mixed batch. PaTAS separates MNIST from Fashion-MNIST at $0.975$ AUROC with a false-positive rate at $95\%$ TPR of $0.145$, against $0.883/0.413$ for the softmax baseline and $0.936/0.310$ for MC~Dropout, and detects corrupted inputs at $0.972$ (MNIST) and $0.835$ (Fashion-MNIST) where every confidence-based baseline remains near chance ($0.51$--$0.63$). The ablation row attributes the mechanism: replacing the conformity-based input opinions with the constant fully-trusted opinion collapses PaTAS to baseline level ($0.896/0.415$ OOD, $0.639$ corruption detection), confirming that the input-level trust of \cref{def:trustfunc}, not the confidence signal, carries the detection.

\begin{table}[ht]
\centering
\setlength{\tabcolsep}{3.5pt}
\begin{tabular}{lcccc}
\toprule
 & \multicolumn{2}{c}{OOD: MNIST vs.\ F-MNIST} & \multicolumn{2}{c}{Corruption det.\ $\uparrow$} \\
Method & AUROC $\uparrow$ & FPR@95 $\downarrow$ & MNIST & F-MNIST \\
\midrule
PaTAS (ours) & $\mathbf{0.975}{\scriptstyle\pm.001}$ & $\mathbf{0.145}$ & $\mathbf{0.972}{\scriptstyle\pm.002}$ & $\mathbf{0.835}{\scriptstyle\pm.003}$ \\
\;\;w/o input trust & $0.896$ & $0.415$ & $0.639$ & -- \\
Softmax & $0.883{\scriptstyle\pm.004}$ & $0.413$ & $0.631$ & $0.558$ \\
MC Dropout & $0.936{\scriptstyle\pm.002}$ & $0.310$ & $0.621$ & $0.545$ \\
EDL & $0.883{\scriptstyle\pm.004}$ & $0.658$ & $0.547$ & $0.510$ \\
\bottomrule
\end{tabular}
\caption{Detection of distribution violations (mean $\pm$ std over three evaluation seeds). OOD: each method's own score separates clean in-distribution test inputs from OOD inputs passed through the identical preprocessing. Corruption detection: AUROC separating clean from corrupted inputs in the mixed batch. The ablation row replaces conformity-based input opinions with the constant fully-trusted opinion.}
\label{tab:comparative}
\end{table}

On in-distribution data the two signals coincide: PaTAS matches the softmax filter on correct-vs-incorrect discrimination on clean test data ($0.979 \pm 0.003$ for both, MNIST), and after identical post-hoc calibration their calibration error is equal ($\mathrm{ECE}=0.006$). Under corruption the signals separate in the expected direction: the corruption barely reduces accuracy ($97.9\% \to 97.6\%$ on the mixed batch), so the softmax retains a higher correct-vs-incorrect AUROC ($0.974$ vs.\ $0.895$) precisely because it ignores the distribution violation that PaTAS flags. Trust and confidence are thus complementary: confidence ranks correctness \emph{within} the training distribution, while trust exposes inputs that \emph{leave} it, the case where confident predictions are unwarranted even when they happen to be correct. Directionality follows from the same reasoning: in the reverse OOD task (Fashion-MNIST in-distribution, MNIST as OOD), MNIST digits occupy a quiet subset of Fashion-MNIST's marginals rather than violating them, and PaTAS ($0.781$) offers no advantage over MC~Dropout ($0.834$); conformity-based trust detects support \emph{violations}, not support \emph{subsets}.

\begin{table}[ht]
\centering
\setlength{\tabcolsep}{4.5pt}
\begin{tabular}{lccccc}
\toprule
Dataset & Med.\ $\sigma_{\mathrm{eff}}$ & $w_{\max}$ & Top-10\% & $w{\geq}4$ & Corr.\ det. \\
\midrule
MNIST & 0.47 & 52.3 & 0.32 & 0.40 & 0.972 \\
Fashion-MNIST & 0.90 & 64.0 & 0.76 & 0.18 & 0.835 \\
GTSRB & 0.98 & 1.8 & 0.14 & 0.00 & 0.545 \\
CIFAR-10 (gray) & 0.23 & 1.1 & 0.11 & 0.00 & -- \\
\bottomrule
\end{tabular}
\caption{A-priori applicability diagnostics: statistics of the conformity model fitted on each dataset's clean training features \emph{before any model is trained} (median effective deviation, maximum reliability weight, evidence share of the top decile, fraction of strong witnesses $w \geq 4$; all quantities defined in \cref{app:conformity}), against the corruption detection subsequently observed. Registered domains concentrate evidence in reliable witnesses; unregistered imagery has none, and input-level conformity is predicted to be uninformative, and for GTSRB this prediction is confirmed.}
\label{tab:diagnostics}
\end{table}

\Cref{tab:diagnostics} explains when this input-level signal exists. Fitting the conformity model on each dataset's training features yields a self-diagnostic: spatially registered domains (MNIST, Fashion-MNIST) concentrate evidence in a minority of tight-marginal ``reliable witness'' features, and the fraction of strong witnesses even predicts the ordering of the observed detection margins ($0.40 \to 0.972$, $0.18 \to 0.835$, $0.00 \to$ chance). Unregistered imagery (GTSRB, CIFAR-10) has uniformly wide marginals, so no feature is a reliable witness and marginal conformity is predicted to be uninformative, which the full GTSRB battery confirms empirically (corruption detection $0.545$; OOD vs.\ grayscale CIFAR-10 $0.661$ vs.\ $0.669$ for softmax). The applicability of input-level trust can therefore be established per dataset before deployment, and richer input-trust models for unregistered domains are an explicit direction for future work.

Finally, the model-side trust provides a training-quality audit that requires \emph{no test data at all}. Training the MNIST network under increasing label-flip rates $p$ and reading the PaTAS feedforward trust under a fully trusted input (\cref{tab:audit}) shows a tripwire behavior: trust drops by $31.7\%$ relative to the clean-trained model at \emph{any} non-zero corruption rate, including $p{=}0.1$ where test accuracy moves by only $3$ points and mean confidence by $14\%$. The flat trust level across rates is the designed saturation of the evidence counting: already at $p{=}0.1$, a 128-sample mini-batch contains corrupted labels with probability $1-0.9^{128} \approx 1-10^{-6}$, so every batch delivers the persistent large-gradient corrections that \textsc{NodeTrust} counts as negative evidence, and the converged trust registers \emph{that} corruption is present rather than how much. Mean softmax confidence, in contrast, tracks the flip rate almost linearly (cross-entropy on $p$-flipped labels converges toward the label purity) and therefore \emph{measures} severity; reading it, however, requires a clean labelled test set, and a moderate confidence value is indistinguishable from a merely difficult task. The two signals answer different questions: PaTAS trust detects \emph{that} training data was corrupted from the training dynamics alone; confidence quantifies the damage given labelled data.

\begin{table}[ht]
\centering
\setlength{\tabcolsep}{5pt}
\begin{tabular}{cccccc}
\toprule
$p$ & Test acc.\ (\%) & Mean conf. & PaTAS trust & $\Delta$trust & $\Delta$conf. \\
\midrule
0   & 97.9 & 0.981 & 0.998 & --     & --     \\
0.1 & 94.7 & 0.841 & 0.682 & 31.7\% & 14.2\% \\
0.2 & 95.7 & 0.762 & 0.682 & 31.7\% & 22.3\% \\
0.3 & 93.7 & 0.684 & 0.682 & 31.7\% & 30.3\% \\
0.4 & 91.2 & 0.589 & 0.682 & 31.7\% & 39.9\% \\
0.5 & 82.8 & 0.476 & 0.681 & 31.7\% & 51.5\% \\
\bottomrule
\end{tabular}
\caption{Training-quality audit (MNIST, 784-512-10): models trained under label-flip rate $p$. PaTAS trust is the feedforward opinion under a fully trusted input, computed from training dynamics alone without test data, and acts as a corruption tripwire ($\Delta$ = relative drop vs.\ $p{=}0$); mean softmax confidence tracks severity but requires a clean labelled test set.}
\label{tab:audit}
\end{table}
